\section{Proofs}\label{app:proofs}
\subsection{Stopped verification}
\begin{proof}[Proof of Proposition~\ref{prop:verification}]
Apply stopped It\^o's formula to $e^{-.04(s-t)}v(s,u_s,x_s)$, localizing first if needed. On the bounded stopped domain the stochastic integral has zero expectation under the stated admissibility and smoothness conditions. Thus
\[
 J_k^a-v(t,u,x)=\E\!\int_t^\tau e^{-.04(s-t)}\cR_k^av(s,u_s,x_s)\,ds
 +\E[e^{-.04(\tau-t)}(g-v)(\tau,u_\tau,x_\tau)].
\]
For every $a$, $\cR_k^av\leq e+q$, giving $J_k^a\leq v+b+(e+q)C_t$, where $C_t=(1-e^{-.04(1-t)})/.04$. For $\pi$, $\cR_k^\pi v\geq-e$, giving $J_k^\pi\geq v-b-eC_t$. Take the supremum over $a$ and subtract. One-sided test functions follow from the same identity; deterministic time envelopes may be integrated to $1$ because $\tau\leq1$.
\end{proof}

\subsection{Improvement, oracle precision, and polling}
\begin{proof}[Proof of Proposition~\ref{prop:monotone}]
Validity of the two enclosures gives $F(y')\geq L(y')>U(y)\geq F(y)$ on acceptance. A rejection restores the original policy and optimizer state. With forcing term $\sigma h^2$, each accepted payoff gain is strictly larger than that quantity. The sum of these gains is bounded by $M-F(y_0)$, yielding the stated conservative integer bound. No assertion about gradient convergence is needed.
\end{proof}
\begin{proof}[Proof of Lemma~\ref{lem:separation}]
For a valid enclosure of width at most $w_j$, its lower endpoint is no less than the true value minus $w_j$, and its upper endpoint is no greater than that value plus $w_j$. Consequently the lower and upper difference endpoints differ from the true payoff difference by at most $2w_j$ in the relevant direction. If $2w_j<\gamma$, the whole difference interval lies strictly on the correct side of the threshold. The displayed refinement count ensures this strict inequality. If the margin is zero, no positive finite width need separate it.
\end{proof}
\begin{proof}[Proof of Theorem~\ref{thm:poll}]
At an unsuccessful poll, $L(y+hd)\leq U(y)+\sigma h^2$ for every $d\in\cD$. Interval widths bounded by $\kappa h^2$ imply
\[
 F(y+hd)-F(y)\leq(\sigma+2\kappa)h^2.
\]
The Lipschitz-gradient inequality on the feasible segment gives
\[
 F(y+hd)-F(y)\geq h\nabla F(y)'d-\frac L2h^2.
\]
Combining these inequalities, maximizing over $d$, and using the cosine measure yields \eqref{eq:stationarity}. Each successful poll increases payoff by more than $\sigma h^2$. The payoff upper bound therefore allows at most the number of successes in Proposition~\ref{prop:monotone} before an unsuccessful complete poll. At most $|\cD|+1$ oracle evaluations per poll, including a conservative incumbent reevaluation, gives \eqref{eq:pollwork}. Multiplying by the uniform oracle cost gives the work bound. The chosen $h$ gives the stationarity tolerance. All feasibility, regularity, and oracle requirements are precisely those in Assumption~\ref{ass:poll}; failure of any requirement does not imply stationarity.
\end{proof}

\subsection{Transport and financing perturbations}
For completeness, the first-moment recursion $m_n=\beta_1m_{n-1}+(1-\beta_1)J_n'g_n$ with zero initialization expands to the coefficients $a_{nr}$ after bias correction. Taylor expansion of $f$ along $-\eta_nD_n^{-1}\widehat m_n$ yields \eqref{eq:transport}; integrating the Lipschitz Jacobian difference along the segment bounds the remainder by $H\norm{\delta\vartheta_n}^2/2$. The executed carrier uses the corresponding sequence of pullbacks and pushforwards but advances economic coordinates linearly, as stated in \eqref{eq:moving}.

\begin{proof}[Proof of Proposition~\ref{prop:gradient}]
The mean-value inequality and the lower bound on $P_a$ imply $|a-a_N|\leq e_P/d_*$. Implicit differentiation gives $a_y=-P_y/P_a$ and $a_{N,y}=-P_{N,y}/P_{N,a}$. Hence the two reduced gradients are $H_y-H_aP_y/P_a$ and $H_{N,y}-H_{N,a}P_{N,y}/P_{N,a}$. Add and subtract products one factor at a time. The payoff-derivative and price-derivative differences contribute $(E_BM_C+M_BE_C)/d_*$; the reciprocal difference contributes at most $M_BM_CE_D/d_*^2$, since both denominators are at least $d_*$. The direct derivative difference contributes $E_A$. The root-shift allowances are already included in the $E$ terms, so no offset error is counted twice.
\end{proof}

\subsection{The stochastic reference and its interval certificates}
\begin{proof}[Proof of Theorem~\ref{thm:reference}]
The controlled residual of $v$ in the reference economy is
\[
 v_t+(.02x+a+.04px)v_x+\frac{.25^2p^2x^2}{2}v_{xx}
 +\frac{.3^2}{2}v_{zz}+b-a^2/2-.04v.
\]
Complete the square in $a$ and $p$, using $v_x=A(1/x+.1z)$ and $v_{xx}=-A/x^2$. With \eqref{eq:transfer}, the result is exactly
\[
 -\tfrac12\{(a-a^*)^2+.0625A(p-p^*)^2\}.
\]
On the closed box, $.32\leq a^*\leq2.1$ and $.512\leq p^*\leq.768$, so the unconstrained maximizers are admissible. The optimal controls are smooth on a neighborhood of the stopped domain. Stopped It\^o's formula applies, with settlement equal to $v$ on every stopping face. The residual is nonpositive for every admissible policy and zero for the displayed optimizer. Thus $V=v$ and the same formula with an arbitrary policy gives \eqref{eq:reference_regret}. Bounding the nonnegative integrand and using $\tau\leq1$ proves \eqref{eq:reference_certificate} for every starting point and time.
\end{proof}
\begin{proof}[Proof of Theorem~\ref{thm:jets}]
Write a point as $m+h$, with $|h_i|\leq r_i$. The second-order integral Taylor formula gives
\[
 e_j(m+h)=e_j(m)+De_j(m)h+\int_0^1(1-s)h'D^2e_j(m+sh)h\,ds.
\]
The integral weight is $1/2$. Absolute values and the symmetric Hessian expansion yield the diagonal half-weights and the single mixed terms in \eqref{eq:jet}. Whole-cell interval Hessians cover every point on the segment. Both first- and second-order bounds are individually valid, so their componentwise minimum is valid. The Hamiltonian deficit is nondecreasing in the nonnegative squared-error bounds and in the positive coefficient $A$; a complete cover and \eqref{eq:reference_certificate} then give the uniform result.
\end{proof}
\begin{proof}[Proof of Lemma~\ref{lem:bernstein}]
The Bernstein basis is nonnegative on $[0,1]$ and sums to one. Hence the polynomial lies between the smallest and largest exact coefficients $\beta_i$. The denominator $x=(5+3\xi)/4$ is positive and bounded below by $(5+3a)/4$ on the cell, proving \eqref{eq:bernstein}. The exact action error before projection is $A(P-1/x)$; projection onto the feasible interval cannot increase its distance from $a^*$. As $A\leq1$ and the portfolio action is exact, the deficit is bounded by $E^2/2$. Apply Theorem~\ref{thm:reference}. The conversion itself uses the identity that a local power coefficient $q_k$ contributes $q_k{ i\choose k}/{n\choose k}$ to Bernstein coefficient $i$, for $k\leq i$; every operation in this conversion is rational.
\end{proof}
